\chapter{Revisão Bibliográfica}
\label{capitulo:revisao_bibliografica}

Este capítulo possui um conteúdo mais dinâmico e dependente do tema proposto no trabalho.
Nele, o autor deverá pesquisar trabalhos correlatos, seja artigos, monografias, dissertações, teses ou até mesmo sistemas, aplicativos ou sites públicos com o mesmo tema.

Entenda que se o tema do trabalho é um sistema acadêmico, seria de bom tom apresentar os principais sistemas acadêmicos do mercado.

Note que esta apresentação poderia ser tendenciosa no sentido em apresentar a função que está querendo melhorar nestes sistemas. 
Lembrando que o que se propõe melhorar foi apresentado no capítulo introdutório na seção \ref{secao:justificativa}.

Este texto introdutório deve apresentar como este capítulo foi organizado apresentando todas as seções e o conteúdo que será explanado em cada uma delas.

Note que além dos trabalhos correlatos (seção \ref{secao:trabalhos_correlatos}), este capítulo também traz a fundamentação teórica (seção \ref{secao:fundamentacao_teorica}).

\section{Fundamentação teórica}
\label{secao:fundamentacao_teorica}

Segundo \citeonline{RaulSidneiWazlawick2021} a fundamentação teórica são conceitos que você obtém na literatura, ou seja, livros, artigos, monografias, dissertações e teses sobre o assunto que está sendo trabalhado.

Um bom exemplo de elementos da fundamentação teórica são as linguagens e os os \textit{frameworks} que foram utilizados.

\subsection{Unified Modeling Language (UML)}

\subsection{Linguagem A}

\subsection{Framework A}

\subsection{Framework B}

\section{Trabalhos correlatos}
\label{secao:trabalhos_correlatos}

Esta seção apresentará o principais trabalhos relacionados ao tema deste estudo.
Devido a vastidão de trabalhos na área, apenas os trabalhos com maior número de citações no ambiente XYZ forma incluídos no estudo.

Aqui pode-se aplicar a técnica de \textbf{revisão sistemática} ou em casos de desenvolvimento de sistemas você pode fazer uma pesquisa de mercado, região ou na Internet dos principais \textit{softwares} sobre o assunto abordado.

Nota-se que para termos técnicos, \textbf{internet} é diferente de \textbf{Internet}. O primeiro termo, internet, refere-se a uma inter-rede e o segundo, Internet, refere-se a rede mundial de computadores.

\begin{citacao}
Um conjunto de redes interconectadas forma uma rede interligada, ou internet. Esses termos serão usados em um sentido genérico, em contraste com a Internet mundial (uma rede interligada específica), que sempre será representada com inicial maiúscula [...]. \cite{Tanenbaum2021}
\end{citacao}

\subsection{Trabalho A}

Apresentação do trabalho A, inclua imagens, algoritmos, protótipos, projetos, diagramas tudo que fizer necessário ao entendimento do que está sendo apontado.

A figura \ref{figura:sistema_correlato_A} apresenta um exemplo de inserção de figuras em trabalhos.
Note toda figura inserida no seu trabalho deve estar referenciada no texto, ou seja, não pode haver uma figura solta sem que o texto esteja a apresentando.
É importante que a apresentação textual venha antes da figura, em alguns casos excepcionais causados pela organização do texto, pode-se apresentar a figura depois do texto.

%\figura{caminho}{tamanho cm}{legenda}{referencia}
\figura{./capitulos/revisao_bibliografica/trabalho_correlato_A.png}{12cm}{Página inicial do site do sistema correlato A.}{figura:sistema_correlato_A}

\subsection{Trabalho B}

Apresentação do trabalho B, inclua imagens, algoritmos, protótipos, projetos, diagramas tudo que fizer necessário ao entendimento do que está sendo apontado.

\subsection{Trabalho C}

Apresentação do trabalho C, inclua imagens, algoritmos, protótipos, projetos, diagramas tudo que fizer necessário ao entendimento do que está sendo apontado.



